%% LyX 2.4.5 created this file.  For more info, see https://www.lyx.org/.
%% Do not edit unless you really know what you are doing.
\documentclass[11pt,english]{article}
\usepackage[T1]{fontenc}
\usepackage[latin9]{inputenc}
\usepackage{babel}
\usepackage{verbatim}
\usepackage{float}
\usepackage{amsmath}
\usepackage{amssymb}
\usepackage{graphicx}
\usepackage{geometry}
\geometry{verbose,tmargin=1.1in,bmargin=1.1in,lmargin=1.1in,rmargin=1.1in}
\usepackage{setspace}
\usepackage[authoryear,round]{natbib}
\onehalfspacing
\usepackage[bookmarks=false,
 breaklinks=false,pdfborder={0 0 1},backref=section,colorlinks=false]
 {hyperref}

\makeatletter

%%%%%%%%%%%%%%%%%%%%%%%%%%%%%% LyX specific LaTeX commands.
%% Because html converters don't know tabularnewline
\providecommand{\tabularnewline}{\\}

%%%%%%%%%%%%%%%%%%%%%%%%%%%%%% User specified LaTeX commands.
\usepackage{babel}
%
\usepackage{amsfonts}
\newtheorem{prop}{Proposition}
\newtheorem{cor}{Corollary}
\newtheorem{lem}{Lemma}

\setcounter{MaxMatrixCols}{30}
\providecommand{\U}[1]{\protect\rule{.1in}{.1in}}
\newenvironment{proof}[1][Proof]{\noindent\textbf{#1.} }{\ \rule{0.5em}{0.5em}}

\newtheorem{theorem}{Theorem}
\newtheorem{acknowledgement}{Acknowledgement}
\newtheorem{algorithm}[theorem]{Algorithm}
\newtheorem{axiom}[theorem]{Axiom}
\newtheorem{assumption}{Assumption}
\newtheorem{case}[theorem]{Case}
\newtheorem{claim}[theorem]{Claim}
\newtheorem{conclusion}[theorem]{Conclusion}
\newtheorem{condition}[theorem]{Condition}
\newtheorem{conjecture}[theorem]{Conjecture}
\newtheorem{corollary}[theorem]{Corollary}
\newtheorem{criterion}[theorem]{Criterion}
\newtheorem{definition}{Definition}
\newtheorem{example}[theorem]{Example}
\newtheorem{exercise}[theorem]{Exercise}
\newtheorem{lemma}{Lemma}
\newtheorem{notation}[theorem]{Notation}
\newtheorem{problem}[theorem]{Problem}
\newtheorem{proposition}{Proposition}
\newtheorem{remark}[theorem]{Remark}
\newtheorem{solution}[theorem]{Solution}
\newtheorem{summary}[theorem]{Summary}
\usepackage{lscape}
\usepackage{placeins}
\usepackage{microtype}
\usepackage{tikz}
\usepackage{pgfplots}
\pgfplotsset{compat=1.18}
\setlength{\emergencystretch}{2.5em}
\sloppy

\makeatother

\begin{document}
\title{Necessity Entrepreneurship: Self-Employment\thanks{We thank}}
\author{David Chivers\thanks{Department of Economics, Durham Business School, DH1 3LB, \protect\protect\href{mailto:david.chivers@durham.ac.uk}{david.chivers@durham.ac.uk}\textbf{ }}
\and Zhigang Feng\thanks{Department of Economics, University of Nebraska, Omaha, NE 68106,
United States, \protect\protect\href{mailto:z.feng2@gmail.com}{z.feng2@gmail.com}.} \and Elisa Keller\thanks{Department of Economics, University of Exeter, \protect\protect\href{mailto:elisa.keller@exeter.ac.uk}{elisa.keller@exeter.ac.uk}}
\and Bo Li\thanks{Department of Economics, Peking University, \protect\protect\href{mailto:lb.harbin@gmail.com\%5C\%2520}{lb.harbin@gmail.com }.} }
\maketitle
\begin{abstract}
TO BE COMPLETED.

\ 

\textbf{JEL Classification}: E23, I10, O40.

\textbf{Keywords}:Occupational Choice, Entrepreneur, Misallocation, 
\end{abstract}
\pagebreak{}

\section{Introduction}

Why do some households become entrepreneurs only when a viable wage
job is unavailable, and how does that differ from entrepreneurship
chosen under strong outside options? Necessity entrepreneurship is
difficult to identify because the object of interest is counterfactual.
The same observed business spell can reflect comparative advantage,
temporary self-employment during job search, insurance against labor-market
risk, or a response to job loss. These cases are unlikely to have
the same implications for firm scale, hiring, or welfare.

Empirical work usually approaches this problem in two ways. Survey-based
evidence asks entrepreneurs about their motives, while status-based
evidence classifies entry by prior labor-market position. Both are
informative but incomplete. In weighted U.S. GEM APS microdata for
2019-{}-2021, 40.9, 49.6, and 45.7 percent of early-stage entrepreneurs
report that they started a business because jobs are scarce, while
wealth or high-income motives remain common at the same time. The
latest U.S. GEM report points in the same direction \citep{gem_us_report_2024},
and global GEM evidence shows that the opportunity-versus-necessity
mix varies systematically across country groups \citep{kelley_singer_herrington_2016}.
Related work on entrepreneurial aspirations points in the same direction
\citep{chivers_2017}. But motive categories can overlap within the
same entrepreneurial spell, so survey data do not map one-for-one
into the model object.

Status-based classifications face a related problem. \citet{fairlie_fossen_2019}
define opportunity entrepreneurship as entry from employment and necessity
entrepreneurship as entry from non-employment, while business-cycle
evidence shows that unemployment-driven and unincorporated entry behave
differently over the cycle \citep{fairlie_2013} \citep{fossen_2021}.
These objects are useful proxies, but they do not reveal whether a
household would still choose entrepreneurship if a viable employment
relationship were available. That counterfactual is central here.

We therefore define opportunity entrepreneurs as households that choose
entrepreneurship whether they enter the period employed or non-employed,
and necessity entrepreneurs as households that choose entrepreneurship
only when they lack an employment relationship. In the self-employment
version of the model, entrepreneurship also splits into nonemployer
self-employment and employer entrepreneurship, so a household can
be an entrepreneur either as a self-employed nonemployer or as an
employer entrepreneur, allowing the paper to speak not only to entry
but also to the composition and scale of entrepreneurial activity.
The framework combines heterogeneous entrepreneurial ability, human
capital, search frictions, entrepreneur closure risk, and social insurance.

This structure matters for interpretation. A downturn or weaker safety
net can shift activity toward low-scale business ownership without
increasing employer-oriented firm creation. Our motivation data therefore
separate zero-employee business counts, business-application flows,
person-level self-employment rates, and motive-based survey evidence
instead of treating them as interchangeable measures of entrepreneurship.
The model is designed to explain those compositional differences rather
than a single aggregate entrepreneurship statistic.

The paper makes three contributions. First, it provides a structural
definition of necessity entrepreneurship that maps directly into state-contingent
occupational choice. Second, it embeds that definition in a self-employment
framework with heterogeneous education, entrepreneurial ability, closure
risk, search frictions, and social insurance, allowing the model to
speak jointly to entry, employer thresholds, and firm scale. Third,
it interprets necessity entrepreneurship as a compositional margin:
weak outside options can push households toward low-scale business
ownership even when employer-oriented entrepreneurship remains weak.

The rest of the paper proceeds as follows. Section 2 presents the
empirical motivation and measurement choices. Section 3 discusses
related literature. Section 4 presents the economic environment. Section
5 characterizes household behavior and equilibrium. Section 6 discusses
calibration. Section 7 presents baseline quantitative results. Section
8 studies policy counterfactuals, and Section 9 concludes.

\section{Empirical motivation and measurement}

We use four public U.S. data objects: business counts by employment
size, business applications, CPS self-employment rates, and GEM motive
evidence. We do not collapse them into a single entrepreneurship measure
because they differ in unit of observation, denominator, and economic
margin. The main empirical motivation in the paper is the combination
of business counts by size and business applications. CPS self-employment
and GEM are treated as companion evidence that clarifies measurement
rather than as headline figures.

Our stock-side series combines Census Nonemployer Statistics and Census
Business Dynamics Statistics to build annual national establishment
counts from 1997 to 2022. The zero-employee series is total nonemployer
establishments, while positive-employment bins use BDS establishment
counts for 1-{}-4, 5-{}-19, and 20+ employees. Over the 2007-{}-2010
Great Recession window, zero-employee establishments rise from 21,708,021
to 22,110,628 (1.85 percent), while 1-{}-4 employee establishments
fall 3.59 percent, 5-{}-19 employee establishments fall 2.57 percent,
and 20+ employee establishments fall 5.82 percent. This object is
therefore best read as evidence on business composition rather than
as a direct measure of startup flows or person-level self-employment.

\begin{figure}[H]
\begin{centering}
\caption{U.S. business counts by employment size, indexed to 1997.}
\par\end{centering}
\centering{}% Auto-generated by pull_us_business_size_cycle_data.py
\begin{tikzpicture}
\begin{axis}[
width=0.82\textwidth,
height=0.50\textwidth,
title={US business counts by employment size},
ylabel={Index (1997=100)},
xlabel={},
xmin=1997, xmax=2022,
ymin=94.4149, ymax=198.6696,
xtick={1997,2002,2007,2012,2017,2022},
xticklabel style={/pgf/number format/fixed},
ymajorgrids=true,
grid style={gray!30},
tick align=outside,
axis line style={black!70},
tick style={black!70},
legend cell align={left},
legend style={draw=none, fill=none, at={(0.02,0.98)}, anchor=north west, font=\small},
]
\path[fill=gray!18,draw=none] (axis cs:2000.5,94.4149) rectangle (axis cs:2001.5,198.6696);
\path[fill=gray!18,draw=none] (axis cs:2007.5,94.4149) rectangle (axis cs:2008.5,198.6696);
\path[fill=gray!18,draw=none] (axis cs:2008.5,94.4149) rectangle (axis cs:2009.5,198.6696);
\path[fill=gray!18,draw=none] (axis cs:2019.5,94.4149) rectangle (axis cs:2020.5,198.6696);
\addplot[very thick, color=blue!65!black] table [x=year, y=establishments_0_index_1997, col sep=comma] {../data/us_business_counts_by_size_annual.csv};
\addplot[very thick, color=green!45!black] table [x=year, y=establishments_1_4_index_1997, col sep=comma] {../data/us_business_counts_by_size_annual.csv};
\addplot[very thick, color=orange!85!black] table [x=year, y=establishments_5_19_index_1997, col sep=comma] {../data/us_business_counts_by_size_annual.csv};
\addplot[very thick, color=red!65!black] table [x=year, y=establishments_20_plus_index_1997, col sep=comma] {../data/us_business_counts_by_size_annual.csv};
\legend{0 employees,1--4 employees,5--19 employees,20+ employees}
\end{axis}
\end{tikzpicture}

\end{figure}

Our flow-side companion uses Census Business Formation Statistics,
accessed through the seasonally adjusted FRED series for total business
applications and high-propensity business applications. Monthly data
are aggregated to annual sums. Over the 2007-{}-2010 window, total
applications fall from 2,652,392 to 2,497,405 (-5.84 percent), while
high-propensity applications fall from 1,493,428 to 1,163,075 (-22.12
percent). The high-propensity share correspondingly declines from
56.30 percent to 46.57 percent. This complements the stock evidence
by showing that employer-oriented business formation weakens even
when zero-employee activity is comparatively resilient.

\begin{figure}[H]
\begin{centering}
\caption{U.S. business applications, indexed to 2005.}
\par\end{centering}
\centering{}% Auto-generated by pull_us_business_applications_companion.py
\begin{tikzpicture}
\begin{groupplot}[
group style={group size=1 by 2, vertical sep=0.95cm},
width=0.82\textwidth,
height=0.28\textwidth,
xmin=2005, xmax=2025,
xtick={2005,2010,2015,2020,2025},
xticklabel style={/pgf/number format/fixed},
ymajorgrids=true,
grid style={gray!30},
tick align=outside,
axis line style={black!70},
tick style={black!70},
]
\nextgroupplot[
ylabel={Index (2005=100)},
xlabel={},
ymin=69.5481, ymax=234.1728,
xticklabel=\empty,
legend cell align={left},
legend style={draw=none, fill=none, at={(0.02,0.98)}, anchor=north west, font=\small},
]
\path[fill=gray!18,draw=none] (axis cs:2007.5,69.5481) rectangle (axis cs:2008.5,234.1728);
\path[fill=gray!18,draw=none] (axis cs:2008.5,69.5481) rectangle (axis cs:2009.5,234.1728);
\path[fill=gray!18,draw=none] (axis cs:2019.5,69.5481) rectangle (axis cs:2020.5,234.1728);
\addplot[very thick, color=black!75] table [x=year, y=applications_total_index_2005, col sep=comma] {../data/us_business_applications_annual.csv};
\addplot[very thick, color=orange!85!black] table [x=year, y=applications_high_propensity_index_2005, col sep=comma] {../data/us_business_applications_annual.csv};
\legend{Total applications,High-propensity applications}
\nextgroupplot[
ylabel={Share (\%)},
xlabel={},
ymin=28.1807, ymax=61.0054,
]
\path[fill=gray!18,draw=none] (axis cs:2007.5,28.1807) rectangle (axis cs:2008.5,61.0054);
\path[fill=gray!18,draw=none] (axis cs:2008.5,28.1807) rectangle (axis cs:2009.5,61.0054);
\path[fill=gray!18,draw=none] (axis cs:2019.5,28.1807) rectangle (axis cs:2020.5,61.0054);
\addplot[very thick, color=teal!60!black, mark=*, mark size=1.7pt] table [x=year, y expr=\thisrow{high_propensity_share}*100, col sep=comma] {../data/us_business_applications_annual.csv};
\end{groupplot}
\end{tikzpicture}

\end{figure}

To avoid treating nonemployer business counts as a person-level self-employment
series, we also build a CPS companion from BLS and FRED monthly levels
for unincorporated self-employed workers, incorporated self-employed
workers, and total employment. We compute annual averages for 2000-{}-2022
and report self-employment rates as shares of total employment. In
the Great Recession window, the total self-employment rate falls from
11.06 percent to 10.70 percent, with both the unincorporated and incorporated
components also declining. This divergence is useful rather than problematic:
an establishment-side nonemployer count and a person-side self-employment
rate capture different margins and should not be combined mechanically
in a single headline statistic.

Finally, GEM provides a motive-based companion rather than a direct
measure of the model object. The public U.S. APS microdata currently
yield a clean weighted TEA sample for 2019-{}-2021, and the official
2023-{}-2024 U.S. GEM report provides a later aggregate update \citep{gem_us_report_2024}.
In the 2019-{}-2021 U.S. microdata, about 41-{}-50 percent of TEA
respondents cite jobs being scarce, 65-{}-73 percent cite a high-income
or wealth motive, and roughly 30-{}-38 percent cite both motives at
once, while only about 8-{}-18 percent cite jobs scarcity alone. The
key lesson is not that GEM solves the measurement problem, but that
necessity-type and opportunity-type signals can coexist within the
same observed entrepreneurial spell. We therefore use GEM to motivate
the paper's identification problem, not to claim that survey motives
map one-for-one into the structural definition of necessity entrepreneurship.

\section{Literature Review}

This paper connects three strands in a single quantitative framework.

First, we build on work that separates opportunity and necessity entrepreneurship
in the data. \citet{fairlie_fossen_2019} propose an operational distinction
based on pre-entry labor market status and show that opportunity and
necessity entrepreneurship have different cyclical patterns. \citet{sedlacek_sterk_2017}
show that startup cohorts differ in growth potential across the business
cycle. This maps directly to our decomposition of entrepreneurial
entry by prior employment state.

Second, we connect to the entrepreneurship-and-frictions literature.
\citet{buera_2009} shows in a dynamic model that borrowing constraints
can materially change both entry and firm scale, with large welfare
effects through undercapitalized businesses. \citet{hurst_pugsley_2011}
documents that many small businesses are not primarily innovation-driven
and that nonpecuniary motives are quantitatively important. This evidence
supports modeling entrepreneurship as a broad occupational margin
rather than only a high-innovation margin.

Third, our labor-market block follows the search-and-matching tradition.
\citet{petrongolo_pissarides_2001} provides the aggregate matching-function
synthesis that motivates our use of separation and job-finding transitions
to discipline occupational flows.

Relative to these strands, our contribution is to embed an explicit
opportunity-versus-necessity distinction in a heterogeneous-agent
model with search frictions, incomplete insurance, and policy counterfactuals.
This structure lets us quantify how unemployment risk and social insurance
jointly shift entrepreneurial entry, firm scale, and composition.

\section{The Model: Economic Environment}

Consider a Lucas (1978) span of control model, where individuals differ
in the level of human capital, $h.$ As in Gomez and Kuhen (2017),
individuals are educated to either the primary $(P)$, secondary $(S)$
or tertiary level $(T)$. The level of education means that a tertiary
educated individuals have a higher draw of human capital, on average,
than secondary educated individuals, who in turn have a higher draw
on average than primary educated individuals. This human capital draw
is positively correlated with managerial ability $x$ as well as labor
productivity $z$.

At the start of the period agents enter as either unemployed $($$i_{w}=0)$,
employed $($$i_{w}=1)$, or an entrepreneur $($$i_{w}=2)$. Individuals
face a choice of whether to become an entrepreneur, to continue working
or to search for work. However, the value of employment is different
depending on whether an agent enters the period as employed $($$i_{w}=1)$
or not. Although there is a risk that an employee may become separated
from a firm with probability $\rho$, this is less risky than searching
for work which has a $p(\theta)$ probability of obtaining employment.
Consequently, the expected value of employment is higher for those
already in employment than those who are not. Additionally, the probability
of separation is lower for individuals with higher levels of human
capital and they are able to command a higher wage due to higher levels
of productivity. So the value of employment, holding other things
equal, is higher for those with higher human capital as well.

The innovation of this model is that the entrepreneurial decision
results in two types of entrepreneurs emerge, opportunity-based entrepreneurs
and necessity-based entrepreneurs. We define opportunity-based entrepreneurs
are those whose value of entrepreneurship is higher, even if they
had employment. Necessity-based entrepreneurs are those who would
not have become entrepreneurs, if they had employment.

\paragraph{Timing: }

The timing of events can be summarized below. Occupational choice
happens at the start of time $t.$ Those choosing to become entrepreneurs
run projects and choose capital and hired labor each period. If hired
labor satisfies $n_{h}=0$ the business is nonemployer self-employment,
while $n_{h}>0$ corresponds to employer entrepreneurship. Entrepreneurial
businesses face closure risk $\rho_{e}(h)$ . In the benchmark, closure
moves the household out of business and into non-employment without
automatic regular unemployment insurance eligibility. Those who were
previously employed within a project can decide to carry on working
but carry a risk of separation with probability $\rho$. If separation
occurs these workers will claim unemployed benefits and be unable
to look for work until the following period. After separation occurs,
vacancies are posted and who were unemployed in the previous period
are matched with firms. Those that are successful and match with with
probability $p(\theta)$ go on to earn a wage. Whereas those that
are unsuccessful will claim unemployment benefits and be unable to
search again until the following period.~

Unlike the search and matching literature there is no wage bargaining
between workers and the firm at the end of matching. Instead the wage
pins down the occupational choice decision which determines those
wanting to become entrepreneurs and those who will search for work.
Crucially, the worker only finds out whether they enter production
and earn a wage after their occupational choice decision.

\begin{figure}[H]
\begin{centering}
\caption{Timing for Matched Workers}
\par\end{centering}
\includegraphics[width=0.82\textwidth]{../figures/enter_period_employed_submission} 
\end{figure}

\begin{figure}[H]
\caption{Timing for Unmatched Workers}

\includegraphics[width=0.82\textwidth]{../figures/enter_period_unemployed_submission} 
\end{figure}


\subsection{Preferences, endowments and technology}

\paragraph{Preferences: }

Consumption by an agent in period $t$ is $c_{t}$, with utility given
by $U(c_{t})$.

\paragraph{Endowments: }

Agents are endowed with an initial capital asset, $a_{0}$, which
can be used as an input in production. Each individual is endowed
with human capital shock $h^{i}$. Human capital $h$ is positively
correlated with managerial talent $x$, labor productivity $z$ and
separation rate $\rho.$

\paragraph{Production: }

Firms use labor $n$ and capital $k$ to produce a single consumption
good, $y$. Capital depreciates at a constant rate of $\delta$.Managers
can operate only one project.The functional form of the production
function is: 
\begin{equation}
y=xk^{\alpha}\left(\ell_{e}+\xi zn_{h}\right)^{\gamma}\ \ \mbox{where}\ \alpha,\gamma>0.\ \label{eq:firm_prod}
\end{equation}
 so self-employment is feasible at $n_{h}=0$ , while employer entrepreneurship
corresponds to $n_{h}>0$ . Managerial talent $x$ is determined by
the Pareto distribution $F(x)=\left[1+\frac{x-\varsigma}{\xi}\right]^{-\vartheta_{i}}$
where $\vartheta$ is the shape parameter, $\varsigma$ the location
parameter and $\xi$ is the scale parameter.

\paragraph{Factor remuneration: }

Firms rent capital at the common market rate $r$.

\paragraph{Government:}

The government runs a balanced budget each period and provides unemployment
insurance which is financed through lump sum taxation $\tau$, which
is shared by the employee $(1-\psi)\tau$ and employer $\psi\tau$.
This program guarantees each household an unemployment benefit of
$b$.

\subsection{The Firm's problem}

The firm's problem is: 
\begin{equation}
\max_{n,k}Xk^{\alpha}n^{\gamma}-(1+\tau\psi)wzn-rk-\kappa\frac{\rho zn}{q(\theta_{t})}\label{eq:max_firm}
\end{equation}
employees must pay a share of unemployment insurance $\tau$ to all
workers who produce output. Firms have to pay search costs $\kappa$
for all vacancies posted. The number of vacancies posted is designed
to replace any lost workers, $\rho zn.$ This, however, is affected
by the firm-hiring rate $q(\theta_{t})=\bar{m}\left(\frac{v_{t}}{u_{t}}\right)^{-\mu}$,
where aggregate vacancies $v_{t}$ is simply the aggregate level of
all those who have separate from firms $\sum\rho zn$ within the period
and $u_{t}$ is the total unemployed looking for work within that
period (see appendix).

When the market is tight and the hiring rate is low, firms have to
search more and thus face higher search costs in order to replace
lost workers. $\text{This is because }0<q(\theta_{t})<1,$so a lower
$q(\theta_{t})$ increases $\frac{\rho zn}{q(\theta_{t})}$. If the
firm lose 10 workers, they want to hire 10 workers to replace them.
But if the hiring rate, $q(\theta_{t}),$ is 50\% then they will face
double the amount of search costs than if $q(\theta_{t})$ were 100\%.
The upshot of this is that when markets are tight, the size of the
desired firm size $n$ decreases ex ante , as the cost of replacing
workers is high and so entrepreneurs would prefer to run smaller firms.

To avoid issues of directed search, we assume that employees are randomly
matched with firms (as in Chivers et al., 2015) and each firm employees
a distribution of workers across the human capital spectrum such that
each firm employs the market average.

\section{Optimal behavior and equilibrium}

In any time period $t$ agents are distinguished by $(h,a,x,i,z)$.
The timing of the economy is given as follows. 
\begin{enumerate}
\item Households enter each new period with assets $a$ and an employment
status $i_{w}$ 
\item Idiosyncratic shocks $h$ is drawn which determines $x$ (managerial
ability) and $z$ (labor productivity), $\rho$ seperation rate conditional
on education level, $P,S$ or $T.$ 
\item Households make an occupation decision: entrepreneur ($\mathcal{I}_{e}=1$)
or worker ($\mathcal{I}_{e}=0$). 
\item Capital and labor markets clear. A fraction $\rho$ of employees separate
from firms and~$p(\theta)$ are hired to replace them 
\item Households choose: consumption ($c$), borrowing/saving ($a'$) 
\end{enumerate}

\subsection{Firm manager}

Firms are distinguished by their productivity realization $x$. In
the self-employment version, entrepreneurs choose capital and hired
labor while also supplying owner labor to the business. Let $n_{h}$
denote hired labor and $\ell_{e}$ denote owner labor. When $n_{h}=0$
the entrepreneur is self-employed; when $n_{h}>0$ the entrepreneur
is an employer. The static problem for a manager with talent $x^{i}$
is  
\begin{equation}
\max_{k,n_{h}\geq0}xk^{\alpha}\left(\ell_{e}+\xi zn_{h}\right)^{\gamma}-zw(1+\psi\tau)n_{h}-rk-\kappa\frac{\rho zn_{h}}{q(\theta)}-f_{hire}\mathbf{1}\{n_{h}>0\}\label{maxP}
\end{equation}

Self-employment is therefore the corner $n_{h}=0$ , while employer
entrepreneurship corresponds to $n_{h}>0$ . For an interior employer
choice, the first-order condition for hired labor is:

\begin{equation}
\gamma xk^{\alpha}\left(\ell_{e}+\xi zn_{h}\right)^{\gamma-1}\xi z=w(1+\psi\tau)+\frac{\rho\kappa}{q(\theta)}\label{maxP-1}
\end{equation}

\begin{equation}
n^{*}_{h}(k,x,\widetilde{w})=\max\left\{ 0,\frac{1}{\xi z}\left[\left(\frac{\gamma xk^{\alpha}\xi z}{w(1+\psi\tau)+\frac{\rho\kappa}{q(\theta)}}\right)^{\frac{1}{1-\gamma}}-\ell_{e}\right]\right\} \label{FOC}
\end{equation}

Substituting (\ref{FOC}) into (\ref{maxP}) yields the entrepreneur's
profit function: 
\begin{equation}
\pi(k,n^{*}_{h},x;\widetilde{w})=xk^{\alpha}\left(\ell_{e}+\xi zn^{*}_{h}\right)^{\gamma}-zn^{*}_{h}w(1+\psi\tau)-rk-\kappa\frac{\rho zn^{*}_{h}}{q(\theta)}-f_{hire}\mathbf{1}\{n^{*}_{h}>0\}\label{profitj}
\end{equation}


\subsubsection{Capital}

Initially we assume capital markets function perfectly, so the interior
capital choice satisfies
\[
\alpha xk^{\alpha-1}\left(\ell_{e}+\xi zn_{h}\right)^{\gamma}=r.
\]


\subsection{Entrepreneurs}

Entrepreneurs maximize expected discounted utility of consumption
\[
\text{with }n_{h}=0\text{ denoting nonemployer self-employment and }n_{h}>0\text{ denoting employer entrepreneurship.}
\]
\[
\max_{\{c_{t},a_{t+1},k_{t},n_{h,t}\}}\mathbb{E}\sum^{\infty}_{t=0}\beta^{t}U(c_{t})
\]
subject to the following budget constraint: 
\begin{align}
c+a' & \leq(1+r-\delta)a+\pi(k,n_{h},x;\tilde{r},w,\tau)-\tau_{\pi}\label{eq:c_worker-3}
\end{align}

\[
a'\geq-\bar{a}
\]


\subsection{Employed}

The workforce (previous periods employed who choose to reapply for
work) maximize expected discounted utility of consumption 
\[
\max_{\{c_{t},a_{t+1}\}}\mathbb{E}\sum^{\infty}_{t=0}\beta^{t}U(c_{t})
\]
subject to the following budget constraint: 
\begin{align}
c+a' & \leq(1+r-\delta)a+(1-\rho)\left(wz(1-(1-\psi)\tau_{y})\right)+\rho b\label{eq:c_worker}
\end{align}

\[
a'\geq-\bar{a}
\]
where $b$ is the unemployment benefit.

\subsection{Unemployed}

The unemployed maximize expected discounted utility of consumption
\[
\max_{\{c_{t},a_{t+1}\}}\mathbb{E}\sum^{\infty}_{t=0}\beta^{t}U(c_{t})
\]
subject to the following budget constraint: 
\begin{align}
c+a' & \leq(1+r-\delta)a+p(\theta)\left(wz(1-(1-\psi)\tau_{y})\right)+(1-p(\theta))b\label{eq:c_worker-2}
\end{align}

\[
a'\geq-\bar{a}
\]
where $b$ is the unemployment benefit.

\subsection{Government }

The government runs a balanced budget with a lump-sum taxes $T(y)=(y-\lambda_{p}y^{1-\tau_{p}})$

\subsection{The household's problem}

Let $i_{w}$ denote the household's labor market status in the previous
period, with $i_{w}=2$ indicating being entrepreneur, $i_{w}=1$
being employed and $i_{w}=0$ being unemployed. Let $\mathcal{I}_{e}$
indicate occupational choice, where $\mathcal{I}_{e}=1$ if the household
is an entrepreneur and $\mathcal{I}_{e}=0$ if the household is a
worker.

Let $\mathbb{E}_{\rho(h)}$ denote the continuation operator for employed
workers who face job separation, let $\mathbb{E}_{p(\theta)}$ denote
the continuation operator for households searching for work, and let
$\mathbb{E}_{\rho_{e}(h)}$ denote the continuation operator for entrepreneurs
facing business closure. We write $\mathbf{V}_{0}$ for the continuation
value after business closure; in the benchmark this state does not
automatically receive regular unemployment insurance. More specifically,
\begin{align*}
\mathbb{E}_{p(\theta)}\mathbf{V}(a',h,x',i'_{w}) & =p(\theta)\mathbf{V}(a',h,x',i'_{w}=1)+(1-p(\theta))\mathbf{V}(a',h,x',i'_{w}=0)
\end{align*}

\begin{align*}
\mathbb{E}_{\rho(h)}\mathbf{V}(a',h,x',i'_{w}) & =(1-\rho(h))\mathbf{V}(a',h,x',i'_{w}=1)+\rho(h)\mathbf{V}(a',h,x',i'_{w}=0)
\end{align*}

\begin{align*}
\mathbb{E}_{\rho_{e}(h)}\mathbf{V}(a',h,x',i'_{w}) & =(1-\rho_{e}(h))\mathbf{V}(a',h,x',i'_{w}=2)+\rho_{e}(h)\mathbf{V}_{0}(a',h,x',i'_{w}=0)
\end{align*}

\begin{align*}
\mathbf{V}(a,h,x,i_{w}) & =\max\left\{ V^{w}(a,h,x,i_{w}),V^{e}(a,h,x)\right\} \\
V^{w}(a,h,x,i_{w}) & =\max_{\left\{ a',c\right\} }\left\{ u(c)+\iota_{i_{w}=1}\beta\mathbb{E}_{\rho(h)}\mathbf{V}(a',h,x',i'_{w})+\iota_{i_{w}\neq1}\beta\mathbb{E}_{p(\theta)}\mathbf{V}(a',h,x',i'_{w})\right\} \\
V^{e}(a,h,x) & =\max_{\left\{ a',c,k,n_{h}\geq0\right\} }\left\{ u(c)+\chi_{se}\mathbf{1}\{n_{h}=0\}+\chi_{emp}\mathbf{1}\{n_{h}>0\}+\beta\mathbb{E}_{\rho_{e}(h)}\mathbf{V}(a',h,x',i'_{w})\right\} 
\end{align*}
subject to

\begin{align}
c+a' & \leq(1-\tilde{{r}}-\delta)a+inc\label{eq:c_worker-1}
\end{align}
\begin{align*}
inc & =\begin{cases}
ra+\widetilde{w}_{n} & \mbox{ if \ensuremath{\mathcal{I}_{e}=0,} \ensuremath{i_{w}=1}}\\
ra+\widetilde{w}_{u} & \mbox{ if \ensuremath{\mathcal{I}_{e}=0,} \ensuremath{i_{w}\neq1}}\\
ra+\pi(k,n_{h},x;\tilde{r},w,\tau) & \mbox{ if \ensuremath{\mathcal{I}_{e}=1}}
\end{cases}
\end{align*}

where $inc$ is the earnings of the worker or entrepreneur, with entrepreneurial
income coming from either self-employment ($n_{h}=0$ ) or employer
entrepreneurship ($n_{h}>0$ ), $\widetilde{w}_{n}=(1-\rho)\left\{ zw(1-(1-\psi)\tau)\right\} +\rho b$
and $\widetilde{w}_{u}=p(\theta)\left\{ zw(1-(1-\psi)\tau)\right\} +\{1-p(\theta)\}b$.

\subsection{Equilibrium}

Let $\hat{h}$ be the critical value of human capital that determines
occupational choice which is affected by both employment status and
education level. If $h\geq\hat{h}$ then individuals choose entrepreneurship
over becoming a worker. The critical value of human capital is always
higher for those who have employment than those who are unmatched
with a firm $\hat{h}_{i_{w}\neq1}<\hat{h}_{i_{w}=1}$. We define necessity
entrepreneurship as those who would not have become an entrepreneur
\textit{if} they had an employment contract: 
\[
\hat{h}_{i_{w}\neq1}<h<\hat{h}_{i_{w}=1}.
\]
We define opportunity-driven entrepreneurship as those who would become
an entrepreneur \textit{regardless} of circumstance: 
\[
h\geq\hat{h}_{i_{w}=1}\geq\hat{h}_{i_{w}\neq1}
\]
As there are three education levels, there are a total of six critical
values of human capital: two for primary educated matched or unmatched,
$\hat{h}^{P}_{i_{w}=1}\text{ and }\hat{h}^{P}_{i_{w}\neq1}$ two for
secondary educated, $\hat{h}^{S}_{i_{w}=1}\text{ and }\hat{h}^{S}_{i_{w}\neq1},$
and two for tertiary educated, $\hat{h}^{T}_{i_{w}=1}\text{ and }\hat{h}^{T}_{i_{w}\neq1}$.

\section{Calibration}
\begin{center}
\begin{table}[htbp]
\caption{Parameter values, baseline economy}

\centering{}\resizebox{\textwidth}{!}{{\small{}%
\begin{tabular}{cccc}
\hline 
{\small Parameters } & {\small Values } & {\small Description } & {\small Comments/observations/target}\tabularnewline
{\small$\rho_{e}$ } & {\small$\{0.168,0.150,0.144\}$ } & {\small Entrepreneur closure rate by $h$ } & {\small BLS survival evidence + Grashuis (2021), scaled to mean 0.15;
see text}\tabularnewline
{\small$f_{hire}$ } & {\small$0$ baseline; $0.25$ in case 147} & {\small Employer-threshold cost} & {\small Initial experimental choice; motivated by first-worker labor-cost
evidence and compliance-cost evidence; see text}\tabularnewline
\hline 
{\small$\beta$ } & {\small$0.94$ } & {\small Discount factor } & \tabularnewline
{\small$\sigma$ } & {\small$2.0$ } & {\small Risk aversion } & {\small$u(c)=\frac{c^{1-\sigma}}{1-\sigma}$}\tabularnewline
\hline 
{\small$\delta$ } & {\small$0.06$ } & {\small Capital depreciation } & \tabularnewline
{\small$\alpha$ } & {\small$0.3207$ } & {\small Capital share } & \tabularnewline
{\small$\gamma$ } & {\small$0.4693$ } & {\small Labour Share } & {\small Frac. of entrepreneurs}\tabularnewline
\hline 
{\small$h$ } & {\small$\{0.109,0.581,0.310\}$ } & {\small Share of population by $h$ } & {\small Barro and Lee (2013)}\tabularnewline
{\small$\vartheta_{e}$ } & {\small$\{0.967,1.007,1.156\}$ } & {\small Shape parameter by $h$ } & {\small Mean firm size, Gomez and Kuhen (2017)}\tabularnewline
{\small$\varsigma$ } & {\small$2.02655$ } & {\small Location parameter } & {\small$F(x)=\left[1+\frac{x-\varsigma}{\xi}\right]^{-\alpha_{i}}$}\tabularnewline
{\small$\xi$ } & {\small$0.1115$ } & {\small Scale parameter } & {\small Gomez and Kuhen (2017)}\tabularnewline
{\small$z$ } & {\small$\{1.0,1.31,1.61\}$ } & {\small Wage premium by $h$ } & {\small Goldin and Katz (2009)}\tabularnewline
{\small$\rho$ } & {\small$\{0.0383,0.0288,0.0136\}$ } & {\small Job separation rate by $h$ } & {\small Cairo and Cajner (2016)}\tabularnewline
\hline 
{\small$\bar{m}$ } & {\small$0.8$ } & {\small Matching efficiency } & {\small Lubik, 2013 $p(\theta)=\bar{m}\left(\frac{v}{u}\right)^{1-\mu}$}\tabularnewline
{\small$\mu$ } & {\small$[0.5,0.7]$ } & {\small Matching elasticity } & {\small Petrongolo and Pissarides (2001)}\tabularnewline
{\small$\kappa$ } & {\small$[0.03,0.045]\cdot w$ } & {\small Search costs } & {\small (Hagedorn and Manovskii, 2008)}\tabularnewline
\hline 
{\small$\tau_{\pi}$ } & {\small$0$ } & {\small Corporation tax } & {\small Set as zero for now}\tabularnewline
{\small$\tau_{y}$ } & {\small$0.151$ } & {\small Income Tax } & {\small$T(y)=(y-\lambda_{p}y^{1-\tau_{p}})$}\tabularnewline
{\small$\lambda_{p}$ } &  & {\small Income Tax } & {\small Determined in equm.}\tabularnewline
{\small$b$ } & {\small$0.40\cdot w$, } & {\small Social security benefits } & {\small Shimmer and Krause and Lubik (2005)}\tabularnewline
{\small$\psi$ } & {\small$0$ } & {\small Share } & \tabularnewline
{\small$\Delta$ } & {\small$0$ } & {\small Credit spread } & \tabularnewline
{\small$\bar{a}$ } & {\small$0$ } & {\small Borrowing constraint } & \tabularnewline
 &  &  & \tabularnewline
\hline 
\end{tabular}}} 
\end{table}
\par\end{center}

In the self-employment extension we introduce an education-specific
entrepreneur closure risk, $\rho_{e}(h)$. Direct U.S. estimates by
education are limited, so the calibration is constructed in two transparent
steps. First, official U.S. business survival evidence from the BLS
and Census implies substantial startup exit, so we set a conservative
average annual entrepreneur closure rate of 0.15, below first-year
startup exit because the model includes incumbent businesses as well
as new entrants. Second, we use the competing-risk estimates in Grashuis
(2021) for transitions from self-employment to unemployment to impose
a modest education gradient. Mapping the model's three education groups
into high school or less, some college, and bachelor's or more, and
normalizing the relative hazards to the model's education shares,
yields $\rho_{e}(h)=\{0.168,0.150,0.144\}$. In the benchmark interpretation,
entrepreneur closure moves the household out of business but does
not automatically confer regular unemployment insurance eligibility,\footnote{This is intended to match regular U.S. unemployment insurance rules
rather than special temporary programs. Regular UI eligibility is
generally tied to prior covered wage employment, whereas purely self-employed
earnings are not automatically covered in the same way. The institutional
benchmark comes from U.S. Department of Labor eligibility guidance
and IRS treatment of unemployment taxes.} so the closure shock disciplines entrepreneurial survival without
mechanically expanding UI receipt.We treat this as an initial calibration
rule rather than a directly estimatedparameter, and use a modest gradient
because the broader literature on educationand business survival is
mixed.

The self-employment branch also allows for an employer-threshold cost,
$f_{hire}$, paid whenever the entrepreneur hires strictly positive
labor. This is separate from the smooth search-and-matching hiring
cost governed by $\kappa$: $\kappa$ captures proportional vacancy
and replacement costs, whereas $f_{hire}$ captures the discrete hurdle
of becoming an employer rather than remaining self-employed. The literature
supports this distinction,\footnote{Recent evidence on the first-worker margin shows that labor costs
can strongly affect the decision to become an employer. Cockx and
Desiere (2024) exploit a Belgian reform that reduced the labor cost
of the first employee by 13 percent and find a 31 percent jump in
first-time employer entry, implying an elasticity of about -2.39.
In related U.S. evidence, Guo and Wallskog (2025) show that even new-employer
payroll taxes averaging about \$350 per worker reduce first-worker
hiring, while Harju, Matikka, and Rauhanen (2019) show that entrepreneurs
respond strongly to compliance costs around tax thresholds. These
papers motivate an employer-threshold cost, but do not pin down a
structural dollar-equivalent fixed cost for our model.} but direct structural calibration of a fixed employer cost remains
limited. We therefore treat $f_{hire}$ as a disciplined experimental
choice rather than a directly estimated parameter. In the current
benchmark with closure risk and partial capital loss we set $f_{hire}=0$,
while the first sensitivity run sets $f_{hire}=0.25$ in model units
to test whether a modest employer hurdle improves the self-employed
versus employer composition.

\section{Baseline Results}

The baseline equilibrium shows moderate entrepreneurship (0.1781 overall),
and the benchmark self-employment calibration places that margin inside
entrepreneurship itself. In case 147, 57.7 percent of entrepreneurs
are self-employed nonemployers and 42.3 percent are employer entrepreneurs.
Throughout the paper, entrepreneur is the broad category, while self-employed
refers to the nonemployer subset of entrepreneurs. The composition
differs sharply by education: low-education entrepreneurs are mostly
self-employed, middle-education entrepreneurs are split more evenly,
and high-education entrepreneurs are predominantly employer entrepreneurs.
Average entrepreneur labor demand is 6.1587 workers, so the key object
in this branch is not a large aggregate entrepreneurship boom but
the composition of entrepreneurial activity across business types
and education groups. The self-employment version of the model adds
a second composition margin within entrepreneurship itself: whether
entrants remain nonemployer self-employed or cross the threshold into
employer entrepreneurship. That margin is central for interpreting
the closure-risk and employer-threshold exercises that follow, even
when changes in aggregate entry are modest.

\begin{table}[H]
\caption{Benchmark self-employment equilibrium: baseline results}
\label{tab:self_employment_benchmark_baseline}
\centering
\small
\resizebox{\textwidth}{!}{%
\begin{tabular}{|l|c|}
\hline
 & Baseline \\
\hline
\multicolumn{2}{|l|}{Entrepreneur share} \\
\hline
All households & 0.1781 \\
\hline
Low education & 0.7187 \\
\hline
Middle education & 0.1329 \\
\hline
High education & 0.0725 \\
\hline
\multicolumn{2}{|l|}{Output per household} \\
\hline
All households & 1288.2102 \\
\hline
Low education & 211.3357 \\
\hline
Middle education & 760.4220 \\
\hline
High education & 2656.0306 \\
\hline
\multicolumn{2}{|l|}{Average entrepreneur size $n$} \\
\hline
All entrepreneurs & 6.1587 \\
\hline
Low education & 1.1767 \\
\hline
Middle education & 6.4030 \\
\hline
High education & 22.6811 \\
\hline
\multicolumn{2}{|l|}{Average employer size $n$} \\
\hline
Employer entrepreneurs, all education & 14.5528 \\
\hline
Low education & 6.7829 \\
\hline
Middle education & 11.8076 \\
\hline
High education & 25.6474 \\
\hline
\multicolumn{2}{|l|}{Average entrepreneur capital $k$} \\
\hline
All entrepreneurs & 98.1596 \\
\hline
Low education & 31.5415 \\
\hline
Middle education & 101.8159 \\
\hline
High education & 317.7547 \\
\hline
\multicolumn{2}{|l|}{Self-employed among entrepreneurs ($n_h = 0$)} \\
\hline
All education & 0.5768 \\
\hline
Low education & 0.8265 \\
\hline
Middle education & 0.4577 \\
\hline
High education & 0.1157 \\
\hline
\multicolumn{2}{|l|}{Employer entrepreneurs among entrepreneurs ($n_h > 0$)} \\
\hline
All education & 0.4232 \\
\hline
Low education & 0.1735 \\
\hline
Middle education & 0.5423 \\
\hline
High education & 0.8843 \\
\hline
\end{tabular}%
}
\end{table}


\section{Quantitative Analysis}

This section evaluates each policy experiment relative to the baseline
and reports how unemployment insurance design and tax closure affect
entry, scale, and output. In the self-employment branch, the key additional
object is the composition of entrepreneurship between nonemployer
self-employment and employer entrepreneurship, not only the overall
entrepreneurship rate.

\begin{table}[H]
\caption{Benchmark self-employment equilibrium: entrepreneur composition by education}
\label{tab:self_employment_benchmark_baseline_vs_noui}
\centering
\small
\begin{tabular}{|l|c|c|}
\hline
 & Baseline (UI=0.40) & No UI (UI=0.00) \\
\hline
Entrepreneur share (all households) & 0.1781 & 0.3425 \\
\hline
Avg entrepreneur labor demand $n$ & 6.1587 & 2.5529 \\
\hline
Avg entrepreneur capital $k$ & 98.1596 & 54.0490 \\
\hline
\multicolumn{3}{|l|}{Self-employed among entrepreneurs ($n_h = 0$)} \\
\hline
All education & 0.5768 & 0.7670 \\
\hline
Low education & 0.8265 & 0.8711 \\
\hline
Middle education & 0.4577 & 0.7732 \\
\hline
High education & 0.1157 & 0.4408 \\
\hline
\multicolumn{3}{|l|}{Employer entrepreneurs among entrepreneurs ($n_h > 0$)} \\
\hline
All education & 0.4232 & 0.2330 \\
\hline
Low education & 0.1735 & 0.1289 \\
\hline
Middle education & 0.5423 & 0.2268 \\
\hline
High education & 0.8843 & 0.5592 \\
\hline
\end{tabular}
\end{table}


\FloatBarrier

\subsection{Low unemployment insurance with endogenous tax}

This subsection reports the low-UI endogenous-tax experiment against
the baseline.

\begin{table}[H]
\caption{Policy experiment (case 101, post-fix): Baseline vs low UI (0.05) with endogenous tax}
\label{tab:exp_lowui_endogtax_case101_postfix}
\centering
\small
\begin{tabular}{|l|c|c|}
\hline
 & Baseline & Low UI (Endogenous Tax) \\
\hline
UI & 0.40 & 0.05 \\
\hline
Interest rate $r$ & 0.0959 & 0.0892 \\
\hline
Wage $w$ & 3.1500 & 3.2511 \\
\hline
Job-finding rate $\theta$ & 0.6111 & 0.6280 \\
\hline
Convergence metric (best\_tol) & 0.0068 & 0.0020 \\
\hline
Equilibrium income tax $\tau_y$ & 0.0181 & 0.0028 \\
\hline
Simulation tax update $\tau_y^{sim}$ & 0.0184 & 0.0021 \\
\hline
Total tax revenue & 6584.28 & 1056.83 \\
\hline
UI outlays (lowerincome) & 8994.34 & 1092.55 \\
\hline
Entrepreneurship share & 0.07884 & 0.07944 \\
\hline
Entrepreneur count (all) & 7884 & 7944 \\
\hline
Entrepreneur count (low edu) & 980 & 975 \\
\hline
Entrepreneur count (med edu) & 4251 & 4424 \\
\hline
Entrepreneur count (high edu) & 2653 & 2545 \\
\hline
Unemployment rate (low edu) & 0.0925 & 0.0857 \\
\hline
Unemployment rate (med edu) & 0.0727 & 0.0686 \\
\hline
Unemployment rate (high edu) & 0.0612 & 0.0533 \\
\hline
Avg entrepreneur labor demand $n$ & 14.9009 & 14.6931 \\
\hline
Mean firm size $n$ (low edu) & 8.9299 & 8.7125 \\
\hline
Mean firm size $n$ (med edu) & 12.0528 & 11.6187 \\
\hline
Mean firm size $n$ (high edu) & 21.5651 & 21.9708 \\
\hline
\end{tabular}
\end{table}


With the post-fix case-101 run at UI=0.05, entrepreneurship increases
only slightly relative to baseline, while tax rates and transfers
fall sharply. Average firm size is close to baseline, with most of
the scale reduction concentrated in low- and medium-education entrepreneurs.

\FloatBarrier

\subsection{No unemployment insurance with endogenous tax}

This subsection reports the no-UI endogenous-tax experiment against
the baseline.

\begin{table}[H]
\caption{Policy experiment (case 101, post-fix): Baseline vs no UI (0.00) with endogenous tax}
\label{tab:exp_noui_endogtax_case101_postfix}
\centering
\small
\begin{tabular}{|l|c|c|}
\hline
 & Baseline & No UI (Endogenous Tax) \\
\hline
UI & 0.40 & 0.00 \\
\hline
Interest rate $r$ & 0.0959 & 0.0849 \\
\hline
Wage $w$ & 3.1500 & 3.3336 \\
\hline
Job-finding rate $\theta$ & 0.6111 & 0.6119 \\
\hline
Convergence metric (best\_tol) & 0.0068 & 0.0056 \\
\hline
Equilibrium income tax $\tau_y$ & 0.0181 & 0.0007 \\
\hline
Simulation tax update $\tau_y^{sim}$ & 0.0184 & 0.0000 \\
\hline
Total tax revenue & 6584.28 & 254.13 \\
\hline
UI outlays (lowerincome) & 8994.34 & 0.00 \\
\hline
Entrepreneurship share & 0.07884 & 0.08466 \\
\hline
Entrepreneur count (all) & 7884 & 8466 \\
\hline
Entrepreneur count (low edu) & 980 & 1081 \\
\hline
Entrepreneur count (med edu) & 4251 & 4773 \\
\hline
Entrepreneur count (high edu) & 2653 & 2612 \\
\hline
Unemployment rate (low edu) & 0.0925 & 0.0908 \\
\hline
Unemployment rate (med edu) & 0.0727 & 0.0733 \\
\hline
Unemployment rate (high edu) & 0.0612 & 0.0576 \\
\hline
Avg entrepreneur labor demand $n$ & 14.9009 & 13.6459 \\
\hline
Mean firm size $n$ (low edu) & 8.9299 & 5.6389 \\
\hline
Mean firm size $n$ (med edu) & 12.0528 & 10.8825 \\
\hline
Mean firm size $n$ (high edu) & 21.5651 & 21.6084 \\
\hline
\end{tabular}
\end{table}


\begin{table}[H]
\caption{Firm size by education (mean entrepreneur labor demand $n$): case 101, post-fix endogenous-tax runs}
\label{tab:firm_size_edu_case101_postfix}
\centering
\small
\begin{tabular}{|l|c|c|c|}
\hline
 & Baseline (UI=0.40) & Low UI (UI=0.05) & No UI (UI=0.00) \\
\hline
Mean $n$ (low edu) & 8.9299 & 8.7125 & 5.6389 \\
\hline
\% vs baseline (low edu) & 0.00\% & -2.43\% & -36.85\% \\
\hline
Mean $n$ (med edu) & 12.0528 & 11.6187 & 10.8825 \\
\hline
\% vs baseline (med edu) & 0.00\% & -3.60\% & -9.71\% \\
\hline
Mean $n$ (high edu) & 21.5651 & 21.9708 & 21.6084 \\
\hline
\% vs baseline (high edu) & 0.00\% & +1.88\% & +0.20\% \\
\hline
Mean $n$ (all entrepreneurs) & 14.9009 & 14.6931 & 13.6459 \\
\hline
\% vs baseline (all entrepreneurs) & 0.00\% & -1.39\% & -8.42\% \\
\hline
Entrepreneurship share & 0.07884 & 0.07944 & 0.08466 \\
\hline
\end{tabular}
\end{table}


Removing UI under endogenous tax raises entry relative to baseline,
but the mechanism is not only a pure risk effect. In this post-fix
run, eliminating UI collapses transfers and pushes the endogenous
tax rate close to zero, while the matching tightness remains nearly
unchanged under the vacancy-floor rule. Firm-size losses are concentrated
in low- and medium-education entrepreneurs, with high-education firm
size nearly unchanged.

\FloatBarrier

\subsection{No unemployment insurance with fixed tax}

This subsection reports the no-UI fixed-tax experiment against the
baseline.

\input{tables/baseline_vs_no_ui_fixed.tex}

With fixed tax and no UI, entry rises modestly, per-capita output
is roughly unchanged, and productivity per entrepreneur declines.

\FloatBarrier

\section{Conclusion}

This paper develops a high-level quantitative framework for studying
necessity entrepreneurship as an equilibrium outcome of occupational
choice under labor market frictions. By defining necessity and opportunity
entrepreneurship through counterfactual choices, the model separates
compositional shifts in entry from changes in aggregate entrepreneurship
rates.

The main message is that unemployment risk and social insurance jointly
shape who becomes an entrepreneur. Reducing insurance can raise entrepreneurship
on paper, but part of this increase reflects necessity entry associated
with lower average firm scale. Conversely, stronger insurance can
reduce necessity entry and improve sorting, even when effects on total
entrepreneurship are modest. These results support evaluating entrepreneurship
policy through composition, not levels alone.

Several extensions remain important for the next draft. First, the
credit friction wedge and job-market experiments should be fully completed
and integrated into a unified welfare comparison. Second, calibration
choices that are still open should be aligned with the code implementation
and reported transparently. Third, the literature discussion can be
broadened further once the project-specific verified PDF library is
restored.

\section*{\newpage}

\section*{Appendix}

\section*{Computation}

Given the values for parameters, the discretized spaces $\Omega_{a}$
for $a$, $\Omega_{z}$ for $z$ and $\Omega_{x}$ for $x$. Each
agent draw a human capital level $h$ (interpreted as educational
attainment) at the beginning of the economy, which will be fixed henceforth.
$h$ will determine their labor productivity $z\in\Omega_{z}$, the
separation rate $\rho$, and managerial ability $x\in\Omega_{x}$.
Let $i_{w}=0$ indicating being entrepreneur, $i_{w}=1$ being employed
and $i_{w}=2$ being unemployed. The numerical algorithm works as
follows. 
\begin{enumerate}
\item Set a tolerance $\epsilon>0$. 
\item Guess $\Phi^{0}=\left(r^{0},w^{0},\rho^{0},\tau^{0}_{y}\right)$.
Solve for optimal household behavior: 
\[
f:\left(\varphi;\Phi\right)\rightarrow\left(c,a',\mathcal{I}_{e},n,k\right),
\]
where $\varphi=\left\{ h,z,x,a,i_{w}\right\} $. We will use the method
of value function iteration as follows.
\begin{enumerate}
\item Guess value function $\mathbf{V}{}^{0}\left(\varphi;\Phi^{0}\right)$
and policy functions $f^{0}\left(\varphi;\Phi^{0}\right)$. 
\item Calculate the cutoff values of $x^{*}_{u}$ (for previously unemployed
worker, or entrepreneur) and $x^{*}_{n}$ (for previously employed
worker) via the following conditions. 
\begin{align}
\widetilde{w}_{u}= & \pi(x^{*}_{u})=\max_{k,n_{h}\geq0}\left\{ x^{*}_{u}k^{\alpha}\left(\ell_{e}+\xi zn_{h}\right)^{\gamma}-w(1+\psi\tau)zn_{h}-rk-\frac{\kappa\rho zn_{h}}{q(\theta)}-f_{hire}\mathbf{1}\{n_{h}>0\}\right\} \label{eq:xstar_u}\\
\widetilde{w}_{n}= & \pi(x^{*}_{n})=\max_{k,n_{h}\geq0}\left\{ x^{*}_{n}k^{\alpha}\left(\ell_{e}+\xi zn_{h}\right)^{\gamma}-w(1+\psi\tau)zn_{h}-rk-\frac{\kappa\rho zn_{h}}{q(\theta)}-f_{hire}\mathbf{1}\{n_{h}>0\}\right\} \label{eq:xstar_n}
\end{align}
where $\widetilde{w}_{n}$ and $\widetilde{w}_{u}$ denote the expected
compensation for the employed worker and unemployed worker respectively.
More specifically, 
\begin{align}
\widetilde{w}_{n} & =(1-\rho)\left\{ zw(1-(1-\psi)\tau)\right\} +\rho b\label{eq:income_n}\\
\widetilde{w}_{u} & =p(\theta)\left\{ zw(1-(1-\psi)\tau)\right\} p(\theta))b.\label{eq:income_u}
\end{align}
Noting that the separation rate $\rho$ is pre-determined by the individual's
initial human capital level $h$. 
\item Update value and policy functions: 
\begin{align}
\mathbf{V}^{1}(h,z,x,a,i_{w})=\label{eq:V_overall}\\
\max_{\left\{ a',c,\mathcal{I}_{e}\right\} }\left\{ \begin{array}{c}
\iota_{i_{w}=0}\cdot\left[\begin{array}{c}
(1-\mathcal{I}_{e})\cdot\mathbb{E}_{p(\theta)}\left(V_{w}+\beta\mathbb{E}\mathbf{V}^{0}(h,z',x',a',i'_{w})\right)\\
+\mathcal{I}_{e}\cdot\left(V_{e}+\beta\mathbb{E}\mathbf{V}^{0}(h,z',x',a',i'_{w}=2)\right)
\end{array}\right]+\\
\iota_{i_{w}=1}\cdot\left[\begin{array}{c}
(1-\mathcal{I}_{e})\cdot\mathbb{E}_{\rho}\left(V_{w}+\beta\mathbb{E}\mathbf{V}^{0}(h,z',x',a',i'_{w})\right)\\
+\mathcal{I}_{e}\cdot\left(V_{e}+\beta\mathbb{E}\mathbf{V}^{0}(h,z',x',a',i'_{w}=2)\right)
\end{array}\right]+\\
\iota_{i_{w}=2}\cdot\left[\begin{array}{c}
(1-\mathcal{I}_{e})\cdot\mathbb{E}_{p(\theta)}\left(V_{w}+\beta\mathbb{E}\mathbf{V}^{0}(h,z',x',a',i'_{w})\right)\\
+\mathcal{I}_{e}\cdot\left(V_{e}+\beta\mathbb{E}\mathbf{V}^{0}(h,z',x',a',i'_{w}=2)\right)
\end{array}\right]
\end{array}\right\} \nonumber \\
f^{1}\left(\varphi;\Phi^{0}\right)=\arg\max\mathbf{V}^{1}\left(\varphi;\Phi^{0}\right)\nonumber 
\end{align}
where $\mathbb{E}_{p(\theta)}\left(V_{w}+\beta\mathbb{E}\mathbf{V}(a',x',i'_{w})\right)$
denotes the expected payoff for the individual who is previously unemployed
worker, or entrepreneur, and decides to look for paid job, while $\mathbb{E}_{\rho}\left(V_{w}+\beta\mathbb{E}\mathbf{V}(a',x',i'_{w})\right)$
represents the expected payoff for the previously employed worker
who decide to stay in workplace. 
\begin{align}
\mathbb{E}_{p(\theta)}\left(V_{w}+\beta\mathbb{E}\mathbf{V}(h,z',x',a',i'_{w})\right) & =p(\theta)\left(V_{n}+\beta\mathbb{E}\mathbf{V}(h,z',x',a',i'_{w}=1)\right)\label{eq:V_u}\\
 & +(1-p(\theta))\left(V_{u}+\beta\mathbb{E}\mathbf{V}(h,z',x',a',i'_{w}=0)\right)\nonumber \\
\mathbb{E}_{\rho}\left(V_{w}+\beta\mathbb{E}\mathbf{V}(h,z',x',a',i'_{w})\right) & =(1-\rho)\left(V_{n}+\beta\mathbb{E}\mathbf{V}(h,z',x',a',i'_{w}=1)\right)\label{eq:V_w}\\
 & +\rho\left(V_{u}+\beta\mathbb{E}\mathbf{V}(h,z',x',a',i'_{w}=0)\right)\nonumber 
\end{align}
To simplify the analysis, we assume that the individual who was previously
employed and newly separated from the firm is not allowed to search
for job at the current period. The measure for the individuals who
are looking for jobs, $u_{t}$, the total vacancies $v_{t}$, and
job matching rate $\theta_{n,t}$ are given as follows. 
\begin{align}
u_{t} & =\left(\int_{x}\iota_{i_{w}=0}+\int_{x}\iota_{i_{w}=2}\right)\cdot\int^{x^{*}_{u}}_{0}\label{eq:measure_u}\\
v_{t} & =\int_{x}\iota_{i_{w}=1}\cdot\rho\label{eq:measure_v}\\
\theta_{n,t} & =\bar{m}\left(\frac{v_{t}}{u_{t}}\right)^{1-\mu}\label{eq:measure_matching}
\end{align}
Note that we can write the above value functions explicitly as a system
of sub-problems: $P^{i_{w}=0}_{x<x^{*}_{u}}$ , $P^{i_{w}=0}_{x\geq x^{*}_{u}}$,
$P^{i_{w}=1}_{x<x^{*}_{n}}$ , $P^{i_{w}=1}_{x\geq x^{*}_{n}}$, $P^{i_{w}=2}_{x<x^{*}_{u}}$
, $P^{i_{w}=2}_{x\geq x^{*}_{u}}$, which correspond to the dynamic
programing problem for i). previously self-employed with $x<x^{*}_{u}$;
ii). previously self-employed with $x\geq x^{*}_{u}$; iii). previously
employed with $x<x^{*}_{n}$; iv). previously employed with $x\geq x^{*}_{n}$;
v). previously unemployed with $x<x^{*}_{u}$; vi). previously unemployed
with $h\geq h^{*}_{u}$. To economize the space, we only provide the
details for $P^{i_{w}=0}_{x<x^{*}_{u}}$ , and $P^{i_{w}=0}_{x\geq x^{*}_{u}}$
as follows. \\
 $\left(P^{i_{w}=0}_{x<x^{*}_{u}}\right)$: 
\begin{equation}
\mathbf{V}^{1}(h,z,x<x^{*}_{u},a,i_{w}=0)=\max_{a'_{n},a'_{u}}\left\{ \begin{array}{c}
p(\theta)\left(u(c_{n})+\beta\mathbb{E}\mathbf{V}^{0}(h,z',x',a'_{n},i'_{w}=1)\right)\\
+(1-p(\theta))\left(u(c_{u})+\beta\mathbb{E}\mathbf{V}^{0}(h,z',x',a'_{u},i'_{w}=0)\right)
\end{array}\right\} \label{eq:V_iw0_xl}
\end{equation}
s.t. 
\begin{align}
c_{n}+a'_{n} & \leq(1-\tilde{r}-\delta)a+zw(1-(1-\psi)\tau),\text{if matched with a firm}\label{eq:bdg_iw0_xl_1}\\
c_{u}+a'_{u} & \leq(1-\tilde{r}-\delta)a+b,\text{if unemployed}\label{eq:bdg_iw0_xl_2}
\end{align}
$\left(P^{i_{w}=0}_{x\geq x^{*}_{u}}\right)$: 
\begin{equation}
\mathbf{V}^{1}(h,z,x\geq x^{*}_{u},a,i_{w}=0)=\max_{a_{e}',k,n_{h}\geq0}\left\{ u(c_{e})+\chi_{se}\mathbf{1}\{n_{h}=0\}+\chi_{emp}\mathbf{1}\{n_{h}>0\}+\beta\left[(1-\rho_{e}(h))\mathbb{E}\mathbf{V}^{0}(h,z',x',a_{e}',i'_{w}=2)+\rho_{e}(h)\mathbb{E}\mathbf{V}^{0}_{0}(h,z',x',a_{e}',i'_{w}=0)\right]\right\} \label{eq:V_iw0_xh}
\end{equation}
s.t. 
\begin{equation}
c_{e}+a_{e}'\leq(1-\tilde{r}-\delta)a+\pi(k,n_{h},x;\tilde{r},w,\tau)\label{eq:bdg_iw0_xh}
\end{equation}
where 
\begin{equation}
\pi(k,n_{h},x;\tilde{r},w,\tau)=xk^{\alpha}\left(\ell_{e}+\xi zn_{h}\right)^{\gamma}-w(1+\psi\tau)zn_{h}-rk-\frac{\kappa\rho zn_{h}}{q(\theta)}-f_{hire}\mathbf{1}\{n_{h}>0\}.\label{eq:profit_emtrepreneur}
\end{equation}
\item Stop if $\max\left\{ \left|\mathbf{V}^{1}-\mathbf{V}^{0}\right|,\left|f^{1}-f^{0}\right|\right\} \leq\epsilon$.
Set $\mathbf{V}^{*}=\mathbf{V}^{1}$, and $f^{*}=f^{1}$. Otherwise,
set $\mathbf{V}^{0}=\mathbf{V}^{1}$, $f^{0}=f^{1}$ and repeat from
step (c). 
\end{enumerate}
\item Generate a large number of individuals, $N=100000$. For each agent
$j$ assign a vector of initial condition $\left(h^{j}_{t},z^{j}_{0},x^{j}_{0},a^{j}_{0},i^{j}_{w,0}\right)$,
where $x^{j}_{0}\sim\Omega_{x}$, $z^{j}_{0}\in\Omega_{z}$, $i^{j}_{w,0}\in\Omega_{w}$. 
\item Simulate the economy for $T$ periods, where $T$ is sufficiently
large. 
\item Calculate the following statistics from the simulated path $\left\{ h^{j}_{t},z^{j}_{t},x^{j}_{t},a^{j}_{t},i^{j}_{w,t},\mathcal{I}^{j}_{e,t},n^{j}_{t},k^{j}_{t}\right\} ^{T}_{t=0}$.
\begin{align}
LS^{0} & =\frac{\sum^{N}_{j=1}\left(\iota_{i_{w}=1}\cdot z^{j}\right)}{N}\label{eq:agg_simulation_1}\\
KS^{0} & =\frac{\sum^{N}_{j=1}\left(\iota_{i_{w}=0}\cdot a^{j}+\iota_{i_{w}=1}\cdot a^{j}+\iota_{i_{w}=2}\cdot a^{j}-\iota_{i_{w}=0}\cdot k^{j}\right)}{\sum^{N}_{j=1}a^{j}}\label{eq:agg_simulation_2}\\
u_{t} & =\left(\int_{h}\iota_{i_{w}=0}+\int_{h}\iota_{i_{w}=2}\right)\cdot\int^{h^{*}_{u}}_{0}\\
v_{t} & =\int_{h}\iota_{i_{w}=1}\cdot\rho\\
p(\theta)_{t} & =\bar{m}\left(\frac{v_{t}}{u_{t}}\right)^{1-\mu}\\
G & =\int\tau wz
\end{align}
and $\tau^{1}_{y}$ that balances the government's budget. 
\item Stop and set $\left(r^{*},w^{*},\theta^{*}\right)=\left(r^{0},w^{0},\theta^{0}\right)$,
if $\max\left\{ LS^{0},KS^{0},\left|\theta^{1}_{n}-\theta^{0}_{n}\right|\right\} \leq\epsilon$.
Otherwise, update aggregate variables (restart from step 2): 
\begin{align*}
r^{0} & =\chi r^{0}+(1-\chi)r^{1}\left(KS^{0}\right)\\
w^{0} & =\chi w^{0}+(1-\chi)w^{1}\left(LS^{0}\right)\\
p(\theta)^{0} & =\chi p(\theta)^{0}+(1-\chi)p(\theta)^{1}\\
\tau^{0}_{y} & =\chi\tau^{0}_{y}+(1-\chi)\tau^{1}_{y}
\end{align*}
where $\chi\in(0,1)$ is the step for updating aggregate variables. 
\end{enumerate}

\section*{}

\section*{Appendix Search and Matching}

I am not sure how much of this is needed for coding purposes but this
is the standard search and matching model. The stock of employed workers
$n_{t}$ can be expressed as the current pool of workers who did not
separate from firms, plus new hires

\[
n_{t}=(1-\rho)n_{t-1}+m_{t-1}
\]
The matching function is Cobb-Douglass where

\[
m_{t}=\bar{m}v^{1-\mu}_{t}u^{\mu}_{t}
\]
$\bar{m}$ is a match efficiency parameter , $v$ is vacancies, $u$
is unemployment, and match elasticity is $0<\mu<1$.The hiring rate
can be defined for firms as total number of new hires to the number
of vacancies

\[
q(\theta_{t})=\frac{m_{t}}{v_{t}}=\bar{m}\theta^{-\mu}_{t}
\]
where $\theta_{t}=\frac{v_{t}}{u_{t}}$is aggregate labour market
tightness. Similarly we can define this from the perspective of the
worker in the job-finding rate.

\[
p(\theta_{t})=\frac{m_{t}}{u_{t}}=\bar{m}\theta^{1-\mu}_{t}
\]
The Job Creation Condition determines the amount of new vacancies
an employer will post. In our model, most workers are on rolling contracts,
so firms only need to post new vacancies after some workers separate
from these rolling contracts. Usually in search and matching models
the wage is determined ex-post of matching by Nash bargaining between
the employer and employee. However, in our span-of-control model the
wage has already been determined ex ante, as the wage is an endogenous
variable that determines occupational choice. The marginal benefit
that the new employee would bring to the firm is expressed as $A$
(which can be derived from the firms production function, $\gamma Xk^{\alpha}(zn)^{\gamma-1}$).
The $JCC\text{ can be then expressed as the costs of replacing out workers on the \ensuremath{LHS} and the expected benefit on the \ensuremath{RHS.}}$

\[
\frac{\kappa}{q(\theta_{t})}=A-w_{t}+\beta E_{t}\left[A-w_{t+1}+(1-\rho)\frac{\kappa}{q(\theta_{t+1})}\right]
\]
where $\kappa$ is the cost of posting a vacancy. %
\begin{comment}
This is different from the standard JCC as usually production starts
the period after matching so the JCC would be $\frac{\kappa}{q(\theta_{t})}=+\beta E_{t}\left[A-w_{t+1}+(1-\rho)\frac{\kappa}{q(\theta_{t+1})}\right].$
See alternative timing version to think about the problem this way,
I am not sure it makes much difference. 
\end{comment}


\subsection*{Matching Steady States}

To find the steady state of the employment equation $n_{t}=(1-\rho)n_{t-1}+m_{t-1}$,
we use the fact that $u=(1-n)$ and $m_{t}=\bar{m}v^{1-\mu}_{t}u^{\mu}_{t}.$
Substituting and getting rid of time subscript gives $(1-u)=(1-\rho)(1-u)+\bar{m}v^{1-\mu}u^{\mu}$
which simplifies to $\bar{m}v^{1-\mu}_{t}u^{\mu}_{t}=$$\rho(1-u)$.
Making $v\text{\,the subject of the foruma gives us the Beveridge curve.}$

\[
v=\left(\frac{\rho(1-u)}{\bar{mu}}\right)^{1-\mu}u
\]
The steady state for the JCC can be found by using $\frac{\kappa}{q(\theta_{t})}=A-w_{t}+\beta E_{t}\left[A-w_{t+1}+(1-\rho)\frac{\kappa}{q(\theta_{t+1})}\right]$
and $q(\theta_{t})=\bar{m}\theta^{-\mu}_{t}=\bar{m}\left(\frac{v_{t}}{u_{t}}\right)^{-\mu}$.
Substituting and getting rid of time subscripts yields 
\[
\frac{\kappa}{\bar{m}}\left(\frac{v}{u}\right)^{\mu}=A-w_{t}+\beta\left[A-w+(1-\rho)\frac{\kappa}{\bar{m}}\left(\frac{v}{u}\right)^{\mu}\right]
\]

Simplifying yields the steady state of the JCC

\[
1-\beta(1-\rho)\frac{\kappa}{\bar{m}}\left(\frac{v}{u}\right)^{\mu}=(1+\beta)\left[A-w\right]
\]

 \bibliographystyle{apalike}
\bibliography{../../../_shared/references/global_references}

\end{document}
